In this manuscript, we focus on derivation of formulas for sums of powers,
by combining
Newton's formula for
powers~\eqref{prop:newtons-interpolation-formula-for-powers}
and hockey stick
identities~\eqref{prop:generalized-hockey-stick-identity},
and~\eqref{lem:shifted-generalized-hockey-stick-identity}.
% to derive formulas for ordinary and multifold sums of powers seamlessly.
This idea, however, can be generalized even further.
Hence, we are free to utilize interpolation formulas for powers
in terms of generalized difference operator $D$,
and binomial coefficients $\binom{f(n)}{k}$.
We are not limited to linear difference operators $\Delta, \nabla$ and $\delta$.
There is a complete theory
on dynamic differential equations on time-scales, founded by
Bohner and Peterson in~\cite{Bohner2001DynamicEO}.
In general, dynamic differential operators assume that the rate of function's
change $h$ is represented by another function $g(h)$,
which is not necessary a linear one.
For example, the Jackson's derivative~\cite{jackson_1909} is as follows.
\begin{align*}
  D_q(f(x)) = \frac{f(qx) - f(x)}{qx-x}
\end{align*}
Thus, finding interpolation formula for $f(x)=x^n$ in terms of $D_q(f(x))$
implies new formulas for sums of powers.
In general, let be an abstract interpolation formula for powers,
\begin{align*}
  n^m = \sum_{k} \binom{f(n)}{k} D(n^m, k).
\end{align*}
Therefore, we obtain one of the formulas for sums of powers,
that involves the abstract difference operator
$D$, evaluated at the point $k$.
By hockey stick identity for binomial coefficients $\binom{n}{k}$,
we get its closed form.
\begin{align*}
  \KnuthRFoldSum{1}{n}{m} = \sum_{k} D(n^m, k) \sum_{j \leq n} \binom{f(j)}{k}.
\end{align*}
Similarly, for multifold sums of powers,
\begin{align*}
  \KnuthRFoldSum{r}{n}{m} = \sum_{k} D(n^m, k) \binom{f(n+r)}{k}.
\end{align*}
Many interpolation approaches involve rising factorials $\risingFactorial{x}{n}$,
or falling factorials $\fallingFactorial{x}{n}$, or regular factorials $n!$,
hence can be expressed
in terms of binomial coefficients, because,
\begin{align*}
  \frac{\fallingFactorial{x}{n}}{n!} = \binom{x}{n}; \quad
  \frac{\risingFactorial{x}{n}}{n!} = \binom{x+n-1}{n}; \quad
  \centralFactorial{x}{n} = n \fallingFactorial{x+\frac{n}{2}-1}{n-1}.
\end{align*}
In particular, Donald Knuth provides the formula for multifold sums of
odd powers in~\cite{knuth1993johann}.
\begin{proposition}[Multifold sums of odd powers]
  For non-negative integers $r,n,m$
  \label{prop:knuth-sums-of-odd-powers}
  \begin{align*}
    \KnuthRFoldSum{r}{n}{2m-1}
    &= \sum_{k=1}^{m} \binom{n+k-1+r}{2k-1+r} \frac{1}{2k} \delta^{2k} 0^{2m} \\
    &= \sum_{k=1}^{m} \binom{n+k-1+r}{2k-1+r} (2k-1)! T(2m, 2k),
  \end{align*}
  where $T(n,k)$ are the central factorial numbers of the second
  kind, see~\cite[section 58]{steffensen1927interpolation}
  and~\cite[formula (10a)]{carlitz1963divided}.
\end{proposition}
This formula originates from Newton's formula in central differences,
evaluated at zero, see~\cite{kolosov_2025_18098442}.
The reason is that central factorial numbers can be expressed in
term of central differences of nothing.
\begin{align*}
  T(n, k) = \frac{1}{k!} \delta^k 0^n = \frac{1}{k!}
  \sum_{j=0}^{k} (-1)^j \binom{k}{j} \left( \frac{k}{2} - j \right)^n
\end{align*}
In general, the
central factorial numbers of the second kind $T(n,k)$ are defined by polynomial
identity~\cite[ch. 6.5, formula (24)]{riordan1968combinatorial}.
Thus,
\begin{lemma} [Riordan power identity]
  For integers $n,m \geq 0$,
  \begin{align*}
    n^m = \sum_{k=1}^{m} T(m,k) \centralFactorial{n}{k}.
  \end{align*}
  where $\centralFactorial{n}{k}$ are central factorials
  $\centralFactorial{n}{k}
  = n \prod_{j=0}^{k-1} \left( n + \frac{k}{2} -j \right)$,
  with $\centralFactorial{n}{0}=1$ for all $n$.
\end{lemma}
The sequence
\href{https://oeis.org/A008957}{\texttt{A008957}}
in the OEIS~\cite{sloane2003line} is the sequence of
central factorial numbers of the second kind $T(2n,2k)$.
As future research direction,
the Knuth's formula~\eqref{prop:knuth-sums-of-odd-powers} utilizes the operator
of central finite differences evaluated in zero.
Thus, it is worth to investigate the existence of sums of odd powers involving
the central differences
evaluated at an arbitrary integer point $t$,
similar to the Theorems~\eqref{theorem:multifold-sums-of-powers-via-newtons-formula},
and~\eqref{theorem:negated-multifold-sums-of-powers-via-newtons-formula}.
